\section{Introduction}
Every semester the timetable committee of the college publishes the class timetable as a multi-page
PDF: one grid per division for S.E., T.E., B.E.\ and M.Tech., with days as rows, time slots as
columns and, below each grid, a legend that expands abbreviations such as \texttt{DBMS}, \texttt{KKD}
or \texttt{LLC}. A single cell may contain a subject, a teacher code, a room number and, for practical
sessions, four stacked lab batches (\texttt{A1}--\texttt{A4}), each with its own teacher and laboratory.

Students and teachers read this document by hand many times a day: to find the next lecture, to see
whether a lab is in 606 or 607, to locate a teacher, or to find an empty room for a make-up lecture.
When a revised PDF is circulated, everyone compares it with the old one by eye, and clashes (one
teacher in two rooms at the same time) are noticed only when they happen. \emph{Dynamic Timetable}
reads the timetable exactly as it is published, checks it, and lets every student and teacher see
their schedule and ask questions about it in plain English on a phone or a laptop.

\subsection{Problem Statement}
\begin{quote}
\emph{Design and implement a system that converts the college's published timetable (PDF, spreadsheet,
Word file or photograph) into verified structured data, keeps every version as an auditable archive,
detects scheduling clashes, and lets each user (student, professor, HOD, Principal or administrator)
get correct answers to everyday timetable questions on a phone or laptop, including questions typed in
plain English, without exposing one user's personal data to another.}
\end{quote}
The difficulties with the current practice are:
\begin{enumerate}
\item \textbf{Unstructured source.} Merged cells, stacked lab rows, page-wise legends, inconsistent
      time labels (the noon row is printed ``12.15 a.m.\ to 01.15 p.m.'') and scanned pages make
      automatic reading hard.
\item \textbf{Manual lookup.} ``Where is my next class?'' requires finding the right page, row,
      column and lab batch and then expanding the legend codes.
\item \textbf{No conflict visibility.} Teacher, room and batch overlaps are not checked mechanically.
\item \textbf{Version confusion.} Old and revised PDFs circulate together.
\item \textbf{Privacy.} A personal timetable (coaching, tuition) must never be visible to others.
\end{enumerate}

\subsection{Objectives}
\begin{enumerate}[label=\textbf{O\arabic*.}]
\item Extract classes from the real timetable PDF with at least 95\% verified automatically, and flag
      every uncertain value for review instead of guessing.
\item Accept PDF, scanned PDF or photo (OCR), Excel/CSV and Word through one pipeline.
\item Provide role-based views for Student, Professor, HOD, Principal and Admin, with permissions
      enforced on the server.
\item Answer everyday questions typed in English with a deterministic engine that is not limited to a
      fixed list of questions, and ask for clarification when a question is ambiguous.
\item Detect professor, room and batch clashes and report only genuine clashes.
\item Keep official and personal timetables strictly separate, with full version history.
\item Answer a question in under 300~ms (median) with no wrong answers under concurrent use.
\item Run on Android, iPhone and laptop browsers from one code base, with a free online demo.
\item Back the system with automated tests, all passing.
\end{enumerate}

\subsection{Scope}
\paragraph{Included.} Upload and extraction (PDF, scanned PDF, JPG/PNG, XLSX/XLS, CSV, DOCX/DOC);
validation, review and audited correction; an archive of all versions with one official timetable;
clash detection; role dashboards; weekly timetable of any batch, teacher or room; natural-language
questions; free-room and teacher-availability search; personal timetables; year-end promotion of
students; administrator master data; audit log.

\paragraph{Not included.} Automatic \emph{generation} of a timetable (the committee still creates it);
attendance; live location of teachers (only the \emph{scheduled} room is reported); exam timetables;
integration with the college ERP.

\paragraph{Users and scale.} Students, professors, HODs, the Principal and the timetable administrator
of one institute: about 2{,}000--3{,}000 users and 500--1{,}000 classes per timetable.

\subsection{Technologies Used}
\begin{table}[H]
\centering
\caption{Technologies used}
\label{tab:tech}
\small
\begin{tabularx}{\textwidth}{@{}p{3.2cm}p{4.6cm}X@{}}
\toprule
\textbf{Layer} & \textbf{Technology} & \textbf{Purpose}\\
\midrule
Back end & Python 3.13, FastAPI~\cite{fastapi} & Typed REST API with automatic OpenAPI documentation\\
Database & PostgreSQL 16~\cite{postgresdocs}, SQLAlchemy 2, Alembic & 31 normalised tables, migrations; also the job queue (\texttt{SKIP LOCKED})\\
Extraction & PyMuPDF~\cite{pymupdf}, Tesseract OCR~\cite{smith2007}, openpyxl, python-docx & Reading PDF grids, scanned pages, Excel and Word files\\
App & React Native~\cite{reactnative}, Expo SDK 58~\cite{expo}, TypeScript, TanStack Query & One code base for Android, iOS and the web\\
Security & Argon2id~\cite{biryukov2016}, JWT~\cite{rfc7519}, RBAC~\cite{sandhu1996} & Password hashing, rotating tokens, role-based access\\
Testing & pytest, Jest, Playwright & Unit, integration and end-to-end tests\\
Deployment & Docker, Render~\cite{render} & Free cloud demo from one blueprint file\\
\bottomrule
\end{tabularx}
\end{table}

\subsection{Assumptions and Constraints}
\paragraph{Assumptions.}
\begin{itemize}
\item The timetable repeats weekly within a semester (date-specific entries are supported but rare).
\item Every user has a college e-mail address issued by the institute.
\item Room floors can be derived from the room number (508 $\rightarrow$ floor 5) until the Admin
      corrects them.
\item A class from 9:00 to 10:00 occupies the interval $[9{:}00, 10{:}00)$, so back-to-back classes do
      not clash.
\end{itemize}
\paragraph{Constraints.}
\begin{itemize}
\item \textbf{Financial:} zero budget, so only open-source software and free hosting were used.
\item \textbf{Technical:} the source is a human-formatted PDF whose layout can change every semester;
      the parser must report what it cannot read.
\item \textbf{Operational:} the official timetable changes only after explicit Admin activation, and
      the Admin can never read personal timetables.
\item \textbf{Time:} one semester alongside regular course work.
\end{itemize}
\newpage
